\subsection{The case of finitely generated modules over a Noetherian ring}

THEOREM 1. Let A be a Noetherian ring and M a finitely generated A-module. There exists a composition series \((\mathbf{M}_i)_{0\leqslant i\leqslant n}\) of M such that, for \(0\leqslant i\leqslant n - 1\) , \(\mathbf{M}_i / \mathbf{M}_{i + 1}\) is isomorphic to \(\mathbf{A} / \mathfrak{p}_i\) , where \(\mathfrak{p}_i\) is a prime ideal of A.

Let \(\mathfrak{E}\) be the set of submodules of M which have a composition series with the property of the statement. As \(\mathfrak{E}\) is non-empty (for \{0\} belongs to \(\mathfrak{E}\) ) and M is Noetherian, \(\mathfrak{E}\) has a maximal element N. If M ≠ N, then M/N ≠ 0 and hence \(\operatorname{Ass}(M/N) \neq \varnothing\) (no. 1, Corollary 1 to Proposition 2); M/N therefore contains a submodule N'/N isomorphic to an A-module of the form A/p, where p is prime; then by definition N' ∈ \(\mathfrak{E}\) , which contradicts the maximal character of N. Then necessarily N = M.

THEOREM 2. Let M be a finitely generated module over a Noetherian ring A and \((\mathbf{M}_{i})_{0\leqslant i\leqslant n}\) a composition series of M such that, for \(0\leqslant i\leqslant n-1\) , \(M_{i}/M_{i+1}\) is isomorphic to \(A/p_{i}\) where \(p_{i}\) is a prime ideal of A. Then

\[
\operatorname{Ass} (\mathbf {M}) \subset \left\{\mathfrak {p} _ {0}, \dots , \mathfrak {p} _ {n - 1} \right\} \subset \operatorname{Supp} (\mathbf {M});\tag{4}
\]

the minimal elements of these three sets are the same and coincide with the minimal elements of the set of prime ideals containing Ann(M).

The inclusion \(\operatorname{Ass}(\mathbf{M}) \subset \{\mathfrak{p}_0, \ldots, \mathfrak{p}_{n-1}\}\) follows immediately from Propositions 1 and 3 of no. 1. For \(0 \leqslant i \leqslant n - 1\) ,

\[
\mathfrak {p} _ {i} \in \operatorname{Supp} (\mathrm{A} / \mathfrak {p} _ {i}) = \operatorname{Supp} (\mathrm{M} _ {i} / \mathrm{M} _ {i + 1})
\]

(Chapter II, § 4, no. 4, Example), whence \(\mathfrak{p}_i \in \operatorname{Supp}(\mathbf{M}_i) \subset \operatorname{Supp}(\mathbf{M})\) (Chapter II, § 4, no. 4, Proposition 16), which shows the inclusion

\[
\left\{\mathfrak {p} _ {0}, \dots , \mathfrak {p} _ {n - 1} \right\} \subset \operatorname{Supp} (\mathbf {M}).
\]

Corollary 1 to Proposition 7 of no. 3 shows that \(\operatorname{Ass}(\mathbf{M})\) and \(\operatorname{Supp}(\mathbf{M})\) have the same minimal elements and (4) shows that these are just the minimal elements of \(\{\mathfrak{p}_0, \ldots, \mathfrak{p}_{n-1}\}\) . The last assertion then follows from Chapter II, § 4, no. 4, Proposition 17.

COROLLARY. If M is a finitely generated module over a Noetherian ring A, Ass(M) is finite.

Under the conditions of Theorem 2, the set \(\{\mathfrak{p}_0,\ldots ,\mathfrak{p}_{n - 1}\}\) is not necessarily determined uniquely by M; in particular it may be distinct from Ass(M) (Exercise 6).

PROPOSITION 8. Let A be a Noetherian ring, a an ideal of A and M a finitely generated A-module. The following conditions are equivalent:

\begin{enumerate}
\def\labelenumi{(\alph{enumi})}
\item
  there exists an element \(x \neq 0\) of \(M\) such that \(ax = 0\) .
\item
  for all \(a \in \mathfrak{a}\) , there exists an element \(x \neq 0\) of \(\mathbf{M}\) such that \(ax = 0\) ;
\item
  there exists \(\mathfrak{p} \in \operatorname{Ass}(\mathbf{M})\) such that \(\mathfrak{a} \subset \mathfrak{p}\) .
\end{enumerate}

Clearly (a) implies (b). By virtue of no. 1, Corollary 2 to Proposition 2, condition (b) means that the ideal a is contained in the union of the prime ideals associated with M and hence in one of them since Ass(M) is finite (Chapter II, § 1, no. 1, Proposition 2); thus (b) implies (c). Finally, if there exists p ∈ Ass(M) such that a ⊂ p, p is the annihilator of an element x ≠ 0 of M (no. 1, Definition 1) and ax = 0; thus (c) implies (a).

PROPOSITION 9. Let A be a Noetherian ring, a an ideal of A and M a finitely generated A-module. For there to exist an integer n \textgreater{} 0 such that \(a^{n}M = 0\) , it is necessary and sufficient that a be contained in the intersection of the prime ideals associated with M.

This intersection is also that of the minimal elements of Supp(M) (no. 3, Corollary 1 to Proposition 7) and to say that a is contained in this intersection is equivalent to saying that V(a) ⊃ Supp(M) in the notation of Chapter II, §4; the conclusion then follows from Chapter II, §4, no. 4, Corollary 2 to Proposition 17.

DEFINITION 2. Given an A-module M, an endomorphism u of M is called almost nilpotent if, for all \(x \in M\) , there exists an integer \(n(x) > 0\) such that \(u^{n(x)}(x) = 0\) .

If M is finitely generated, every almost nilpotent endomorphism is nilpotent.

COROLLARY. Let A be a Noetherian ring. M an A-module and a an element of A. For the homomorphism \(a_{\mathbf{M}}: x \mapsto ax\) of \(\mathbf{M}\) to be almost nilpotent, it is necessary and sufficient that a belong to every ideal of \(\operatorname{Ass}(\mathbf{M})\) .

The condition of the statement is equivalent to saying that for all \(x \in M\) there exists \(n(x) > 0\) such that \((\mathbf{A}a)^{n(x)}(\mathbf{A}x) = 0\) ; by Proposition 9 this means also that a belongs to all the prime ideals associated with the submodule Ax of M; the corollary then follows from the fact that \(\operatorname{Ass}(M)\) is the union of the \(\operatorname{Ass}(Ax)\) where x runs through M (no. 1, formula (1)).

PROPOSITION 10. Let A be a Noetherian ring, E a finitely generated A-module and F an A-module. Then

\[
\operatorname{Ass} \left(\operatorname{Hom} _ {\mathbf {A}} (\mathrm{E}, \mathrm{F})\right) = \operatorname{Ass} (\mathrm{F}) \cap \operatorname{Supp} (\mathrm{E}).\tag{5}
\]

By hypothesis, E is isomorphic to an A-module of the form \(A^{n}/R\) , hence \(\operatorname{Hom}_{\mathbf{A}}(\mathbf{E}, \mathbf{F})\) is isomorphic to a submodule of \(\operatorname{Hom}_{\mathbf{A}}(\mathbf{A}^{n}, \mathbf{F})\) and the latter is isomorphic to \(F^{n}\) ; now, \(\operatorname{Ass}(\mathbf{F}^{n}) = \operatorname{Ass}(\mathbf{F})\) (no. 1, Corollary 1 to Proposition 3); thus \(\operatorname{Ass}(\operatorname{Hom}_{\mathbf{A}}(\mathbf{E}, \mathbf{F})) \subset \operatorname{Ass}(\mathbf{F})\) . On the other hand,

\[
\operatorname{Supp} (\operatorname{Hom} _ {\mathbf {A}} (\mathbf {E}, \mathbf {F})) \subset \operatorname{Supp} (\mathbf {E}):
\]

for every prime ideal p of A, \(\operatorname{Hom}_{\mathbf{A}_{p}}(\mathbf{E}_{p}, \mathbf{F}_{p})\) is isomorphic to \((\operatorname{Hom}_{\mathbf{A}}(\mathbf{E}, \mathbf{F}))_{p}\) (Chapter II, §2, no. 7, Proposition 19), whence our assertion immediately; then we conclude from Theorem 2 that

\[
\operatorname{Ass} \left(\operatorname{Hom} _ {\mathbf {A}} (\mathrm{E}, \mathrm{F})\right) \subset \operatorname{Supp} (\mathrm{E}).
\]

Conversely, let p be a prime ideal of A belonging to \(\operatorname{Ass}(\mathbf{F}) \cap \operatorname{Supp}(\mathbf{E})\) . By definition, F contains a submodule isomorphic to A/p.~On the other hand, since E is finitely generated and \(E_{p} \neq 0\) , we know that there exists a homomorphism \(w \neq 0\) from E to A/p (Chapter II, §4, no. 4, Proposition 20). As there exists an injective homomorphism j from A/p to F, \(j \circ w \in \operatorname{Hom}(\mathbf{E}, \mathbf{F})\) and \(j \circ w \neq 0\) . On the other hand, the relation aw = 0 for some \(a \in A\) is equivalent to \(a \in p\) , the annihilator of every element \(\neq 0\) of A/p being p; then certainly \(p \in \operatorname{Ass}(\operatorname{Hom}_{\mathbf{A}}(\mathbf{E}, \mathbf{F}))\) .

